\section{\ours{}}\label{sec:method}
In the following, we introduce \oursacro{} in two steps.
First, we introduce a non-localized version, termed SVGP-DKL, by combining the SVGP with Deep Kernel Learning (DKL)~\citep{wilson2015Deep}.
In DKL one uses deep neural networks to encode the data in a learned, potentially lower-dimensional, space before applying the GP. Both are then trained end-to-end.
Second, we discuss how we localize the SVGP-DKL to obtain \oursacro{} fulfilling all desiderata.

\textbf{SVGP-DKL}:
To leverage the predictive performance of GNN-based molecular force fields, we use such a GNN as encoder for DKL.
Further, we avoid the expensive runtimes of the GP by relying on the SVGP approximation.
Let $h(\pos, \inputfeatures) = (h_\text{pred} \circ h_{\text{rep}})(\pos, \inputfeatures) $ be a GNN with $h_\text{rep}:\real^{n\times 3} \times \real^{n \times h} \rightarrow \real^{n \times d}$ being its learned invariant atom representations and $h_\text{pred}:\real^{n \times d} \rightarrow \real$ its property predictor.
By replacing the predictor $h_\text{pred}$ with an SVGP, we combine the benefits of GNNs with approximate GPs.
Note that GPs operate on fixed-dimensional spaces and we, thus, have to introduce a summation over all atoms to obtain a representation for the whole molecule:
\begin{equation}
\begin{aligned}
p(E_\star \vert \pos, \inputfeatures) &= \gp_\phi \circ \sum_{i=1}^n h_{\text{rep}}(\pos, \inputfeatures)_i\\ &\sim \normal\left(E_\star \vert \EX [\energy_\star], \Cov(\energy_\star)\right)\label{eq:global_gp}
\end{aligned}
\end{equation}
We can view this as a learned kernel \cite{wilson2015Deep}: 
\begin{equation}
\begin{aligned}
    &k_{\text{SVGP-DKL}}((\matr X^{(k)}, \matr H^{(k)}),(\matr X^{(j)}, \matr H^{(j)})) =\\ k&\left(\sum_i h_{\text{rep}}(\matr X^{(k)}, \matr H^{(k)})_i, \sum_i h_{\text{rep}}(\matr X^{(j)}, \matr H^{(j)})_i\right)
\end{aligned}
\end{equation}
Hence, the prediction of the model is a normal distribution. We can use the variance of the predictive distribution as energy uncertainty estimate. Given that the embeddings of the studied GNNs (SchNet, DimeNet++, NequIP) are invariant to the euclidean group $\set E(3)$ and permutations $S_n$, this GP extension fulfills \desone{}.  

To fulfill the physical desideratum of \destwo{}, we need to calculate the forces as the derivative of the energy. Hence, our method needs to be twice differentiable to be trained. Taking the derivative of the energy with respect to the position thus yields the force prediction: 
\begin{equation}
\matr F = -\frac{\partial E}{\partial \pos} = -\frac{\partial \text{GP}_\phi}{\partial h_{\text{rep}}}\frac{\partial h_{\text{rep}}}{\partial \pos}. 
\end{equation}
As a GP with RBF kernel is $C^\infty$ smooth, we only have to pick a twice differentiable representation function $h_{\text{rep}}$ to fulfill \destwo{}. 
To train the GP, we use the predictive log-likelihood as objective to estimate the energy error \cite{jankowiak_parametric_2020}. This objective, similar to the variational ELBO shown in \autoref{eq:elbo}, can be formally written as: 
\begin{equation}
    \mathcal{L} = \log \E_{q(\vect u)p(\pfun|\vect u)}[p(\pfun_*|\pfun)] - D_{\text{KL}(q(\vect u)||p(\vect u))}.
\end{equation}
This formulation has been reported to yield improved estimates of predictive variances when compared to alternative approaches \cite{jankowiak_parametric_2020, 9982167}. We use the RMSE loss to estimate the force error.
However, as the GP operates on the singular, fixed-size, molecule embedding, it fails to satisfy \desfour{}.

\textbf{Localizing SVGP-DKL}:
To attain \desfour{}, we use a localized version of the just defined GNN-based GP to obtain \ours{} (\oursacro{}). 
We achieve this by reversing the order of the sum and the GP in \autoref{eq:global_gp}.
Instead of using the global graph embedding, we fit the GP on the atomic embeddings directly and our final output is the sum of these atomic GPs. Hence, the predictive distribution of the predicted energy value $\energy_*$ can be written as\footnote{Assuming independence of the atom random variables and, thus, summing only the diagonal elements leads to better performance.}: 
\begin{align}
    p(E_\star \vert \pos, \inputfeatures) &= \sum_i^\natoms \gp_\phi \circ h_{\text{rep}}(\pos, \inputfeatures)_{i}\\  
    &\sim \normal\left(E_\star \,\vert\,\sum_i^\natoms \EX [\energy_\star]_i, \sum_{ij}^\natoms \Cov(\energy_\star)_{ij}\right).\nonumber
\end{align}
Since we sum the individual atom energy contributions and uncertainties, we fulfill \desfour{}. Compared to the global embedding version, adding an independent molecule will increase the uncertainty proportionally. We will show this comparison also empirically in \autoref{sec:experiments}.
In regards to the `application-focused' desiderata, \oursacro{} only needs a single forward pass fulfilling \desseven{}, and as we will show in \autoref{sec:experiments} also fulfills \desfive{} and \dessix{}. 
Thus, it is the only approach fulfilling all of our set desiderata.


\textbf{Avoiding pitfalls.}
As the use of large encoder networks can lead to unstable uncertainty estimation for deep kernel learning models~\citep{ober_promises_2021}, proper regularization has to be deployed.
Ober et al. suggest using a full-bayesian treatment to avoid this pitfall. We use two different approaches: (1) approximate a full bayesian treatment by using dropout \cite{gal2016dropout, srivastava2014dropout} and (2) a fixed-encoder which is motivated by the strong empirical performance in their work (see Table 3 \citet{ober_promises_2021}).

In order to apply dropout to equivariant networks like NequIP, one cannot simply drop any features.
In \autoref{app:baselinedetails}, we discuss how we translate dropout to such group-equivariant neural networks.
Note that during inference, we do not perform dropout with \oursacro{} as we use the variance of the GP's predictive distribution as an uncertainty estimate rather than an empirical Monte Carlo estimate.

Instead of applying dropout, one might also avoid overfitting by fixing the representation function $h_\text{rep}$.
This could be done, e.g., by first training the GNN with its original predictor on the objective.
After pretraining, the representation function will be frozen similar to common fine-tuning on computer vision models.
By limiting the embedding space in such a way we can avoid DKL pitfalls~\citep{ober_promises_2021}.

In the following, we refer to the dropout training approach as \oursacro{} while we call the fixed encoder approach fixed-\oursacro{}.
